\chapter{A Concrete Naming Certificate for $\KernelInductiveAcceptanceUp$}
\label{ch:concrete-instances-kernelinductiveacceptance-namecert}
\origin{ai}

The $\KernelInductiveAcceptanceUp$ packet realizes the inductive-declaration branch of \autoref{ch:visions-kernel-acceptance-as-audit-witness} as finite NameCert data. It sits beside the accepted-declaration packet of \autoref{ch:concrete-instances-kernelacceptancewitness-namecert}, the generic acceptance trace of \autoref{ch:concrete-instances-kernelacceptancetrace-namecert}, the generator closure classifier of \autoref{ch:concrete-instances-generatorclosure-namecert}, the type-classifier membership surface of \autoref{ch:concrete-instances-typeclassifiermembership-namecert}, and the subject-reduction boundary of \autoref{ch:concrete-instances-subjectreduction-namecert}. Those siblings expose acceptance, trace, closure, membership, and preservation rows; this packet isolates the constructor-signature plus emitted-eliminator ledger required for accepted inductive declarations.

\begin{definition}[KernelInductiveAcceptance carrier]
\label{def:kernel-inductive-acceptance-carrier}
A $\KernelInductiveAcceptanceUp$ carrier is a finite $\BHist$ packet
$$
I=(D,S,E,P,R,H,C,Q,N)
$$
with the following displayed rows:
\begin{enumerate}
\item $D$ is the inductive declaration candidate row, including the carrier name, parameters, indices, universe data, and constructor-name inventory;
\item $S$ is the constructor-signature ledger, recording the finite domain and target signature of every constructor accepted for $D$;
\item $E$ is the emitted-eliminator ledger, recording the recursor, induction principle, and computation rows generated from the accepted signature;
\item $P$ is the positivity-refusal row, recording that a candidate whose constructor signature violates the local positivity discipline is refused rather than entered into the environment;
\item $R$ is the recursion-discipline row, recording that eliminator and recursor consumers may read only the emitted rows in $E$;
\item $H$ records componentwise $\hsame$ transports for $D,S,E,P,R,Q,N$;
\item $C$ records the displayed $\Cont$ routes from the declaration candidate to constructor signatures, from constructor signatures to emitted eliminators, from failed signatures to positivity refusal, and from emitted eliminators to downstream consumers;
\item $Q$ is the $\Pkg$ provenance over \autoref{ch:visions-kernel-acceptance-as-audit-witness}, \autoref{ch:concrete-instances-kernelacceptancewitness-namecert}, \autoref{ch:concrete-instances-kernelacceptancetrace-namecert}, \autoref{ch:concrete-instances-generatorclosure-namecert}, \autoref{ch:concrete-instances-typeclassifiermembership-namecert}, \autoref{ch:concrete-instances-subjectreduction-namecert}, $\BHist$, $\hsame$, $\Cont$, and $\NameCert$ rows;
\item $N$ is the local $\NameCert_{\KernelInductiveAcceptanceUp}$ row.
\end{enumerate}
The carrier contains no theorem-proof row, no generic theorem trace, no generator-closure theorem, no type-membership theorem, and no subject-reduction target row beyond the emitted eliminator ledger displayed in $E$.
\end{definition}
\leandef{BEDC.Derived.KernelInductiveAcceptanceUp.KernelInductiveAcceptanceUp}

\begin{theorem}[KernelInductiveAcceptance NameCert obligations]
\label{thm:kernel-inductive-acceptance-namecert-obligations}
For the carrier of \autoref{def:kernel-inductive-acceptance-carrier}, the local $\NameCert_{\KernelInductiveAcceptanceUp}$ surface has five rows:
\begin{enumerate}
\item carrier admission for the finite packet $I=(D,S,E,P,R,H,C,Q,N)$;
\item constructor-signature exactness, requiring every accepted constructor read to be one of the displayed entries of $S$;
\item emitted-eliminator ledger exactness, requiring every recursor, induction-principle, and computation read to pass through $E$ after $S$ is visible;
\item positivity-refusal stability, requiring a failed positivity check to land in $P$ rather than in an environment-entry row;
\item environment-entry non-escape, requiring downstream acceptance, trace, generator, type-classifier, and subject-reduction consumers to read only $D,S,E,P,R,H,C,Q,N$ and not a hidden kernel-correctness theorem.
\end{enumerate}
\end{theorem}
\begin{proof}
Carrier admission is projection from the displayed packet. The declaration candidate, constructor-signature ledger, emitted-eliminator ledger, positivity refusal, recursion discipline, transport, route, provenance, and local name are all coordinates of the same finite $\BHist$ object.

Constructor-signature exactness and emitted-eliminator exactness are the two accepted rows. The accepted constructor inventory is read from $S$, and each eliminator or recursor consumer must first expose $S$ before reading the generated ledger $E$.

The refusal and environment clauses separate accepted and rejected packets. A failed positivity read lands in $P$ and never reaches the environment-entry route, while an accepted read may enter the environment only through the displayed $D,S,E,R,H,C,Q,N$ rows.

The sibling chapters then consume only the finite rows they name. KernelAcceptanceWitness and KernelAcceptanceTrace can read the environment entry, GeneratorClosure can read the accepted generator boundary, TypeClassifierMembership can read emitted type-classifier data, and SubjectReduction can read recursor-facing preservation boundaries; none of those reads adds an unrecorded proof of kernel correctness.
\end{proof}

\begin{theorem}[KernelInductiveAcceptance nontrivial boundary]
\label{thm:kernel-inductive-acceptance-nontrivial-boundary}
The $\KernelInductiveAcceptanceUp$ packet has a nontrivial boundary for TasteGate: an accepted packet displays at least one accepted environment-entry route through $D,S,E,R,H,C,Q,N$, while a candidate whose constructor signature lacks the required positivity row lands in $P$ and is not accepted. Phase C must realize this as a Nontrivial $\KernelInductiveAcceptanceUp$ instance in $\mathsf{BEDC.Derived.KernelInductiveAcceptanceUp.TasteGate}$.
\end{theorem}
\begin{proof}
Use the carrier to separate accepted and refused rows. Accepted packets expose constructor signatures and emitted eliminators through $S$ and $E$, while refused packets expose the positivity boundary $P$.

The NameCert obligation surface makes the separation observable. Environment-entry non-escape permits an accepted route only after the displayed constructor-signature and eliminator ledgers are present, and positivity-refusal stability prevents a failed candidate from taking that route.

Thus the TasteGate boundary is not empty: it has an accepted side and a refused side, both represented by finite rows of the same carrier.
\end{proof}

\begin{theorem}[KernelInductiveAcceptance field-faithful boundary]
\label{thm:kernel-inductive-acceptance-field-faithful-boundary}
The $\KernelInductiveAcceptanceUp$ packet is field-faithful for TasteGate: any accepted read that preserves $D,S,E,P,R,H,C,Q,N$ componentwise preserves the inductive-declaration acceptance row, and any read that drops either the constructor-signature ledger $S$ or the emitted-eliminator ledger $E$ is not a faithful read of this carrier. Phase C must realize this as a FieldFaithful $\KernelInductiveAcceptanceUp$ instance in $\mathsf{BEDC.Derived.KernelInductiveAcceptanceUp.TasteGate}$.
\end{theorem}
\begin{proof}
Start with two accepted reads matched componentwise by the displayed $\hsame$ row. The declaration candidate, signature ledger, eliminator ledger, refusal boundary, recursion discipline, routes, provenance, and local name are therefore transported together.

Apply constructor-signature exactness and emitted-eliminator exactness from the NameCert surface. These two rows are the fields that distinguish inductive acceptance from a generic accepted declaration or trace.

If either row is absent, the carrier no longer exposes the paired signature-plus-eliminator surface. The read may belong to a sibling acceptance packet, but it is not faithful to $\KernelInductiveAcceptanceUp$.
\end{proof}

\begin{theorem}[KernelInductiveAcceptance constructor ledger exhaustion]
\label{thm:kernel-inductive-acceptance-constructor-ledger-exhaustion}
Every accepted constructor read of a $\KernelInductiveAcceptanceUp$ packet is exhausted by the displayed declaration row $D$, constructor-signature ledger $S$, componentwise $\hsame$ transports $H$, displayed $\Cont$ routes $C$, package provenance $Q$, and local $\NameCert$ row $N$. No accepted constructor consumer can read a constructor signature outside the finite ledger $S$.
\end{theorem}
\begin{proof}
Use the carrier of \autoref{def:kernel-inductive-acceptance-carrier}. The declaration candidate $D$, constructor-signature ledger $S$, transports $H$, routes $C$, provenance $Q$, and naming row $N$ are displayed coordinates of the same finite $\BHist$ packet.

Apply the constructor-signature exactness clause of \autoref{thm:kernel-inductive-acceptance-namecert-obligations}. It requires every accepted constructor read to be one of the displayed entries of $S$, with its declaration source and route recorded by the surrounding rows.

The TasteGate commitments are preserved because the read is still through the paired inductive-declaration surface. A consumer that drops $S$ or replaces it by a hidden theorem-proof row leaves the field-faithful boundary of \autoref{thm:kernel-inductive-acceptance-field-faithful-boundary}.
\end{proof}

\begin{theorem}[KernelInductiveAcceptance eliminator ledger exhaustion]
\label{thm:kernel-inductive-acceptance-eliminator-ledger-exhaustion}
Every recursor, induction-principle, or computation read of an accepted $\KernelInductiveAcceptanceUp$ packet factors through the constructor-signature ledger $S$ before entering the emitted-eliminator ledger $E$. No eliminator consumer may read a hidden theorem-proof row or bypass $S$ on the way to $E$.
\end{theorem}
\begin{proof}
Use the carrier of \autoref{def:kernel-inductive-acceptance-carrier}. The rows $D,S,E,R,H,C,Q,N$ are displayed coordinates of the same finite $\BHist$ packet, so a recursor or induction-principle consumer has no carrier coordinate outside them.

Apply the constructor-signature exactness and emitted-eliminator exactness clauses of \autoref{thm:kernel-inductive-acceptance-namecert-obligations}. They force each accepted eliminator read to expose $S$ before $E$.

Finally use the sibling anchors named in the chapter introduction. The accepted-declaration, trace, generator, type-classifier, and subject-reduction chapters consume the emitted route after $E$ is visible, but none of them inserts an extra theorem-proof row between $S$ and $E$.
\end{proof}

\begin{theorem}[KernelInductiveAcceptance refusal nonescape]
\label{thm:kernel-inductive-acceptance-refusal-nonescape}
Every constructor-signature candidate for $\KernelInductiveAcceptanceUp$ that fails the local positivity row lands in the refusal row $P$ and cannot be transported into the accepted environment-entry route $D,S,E,R,H,C,Q,N$. The failed candidate remains outside the accepted declaration, trace, generator, type-classifier, and subject-reduction consumer paths.
\end{theorem}
\begin{proof}
Begin with the carrier row $P$ of \autoref{def:kernel-inductive-acceptance-carrier}. It is the displayed place for failed positivity candidates, distinct from the accepted constructor-signature ledger $S$ and emitted-eliminator ledger $E$.

Apply the positivity-refusal stability clause of \autoref{thm:kernel-inductive-acceptance-namecert-obligations}. A failed signature is routed to $P$ rather than to the accepted environment-entry path.

The nontrivial boundary of \autoref{thm:kernel-inductive-acceptance-nontrivial-boundary} separates the accepted and refused sides of the TasteGate surface. Thus the failed candidate cannot be read by the accepted-declaration, trace, generator, type-classifier, or subject-reduction consumers as an entered inductive declaration.
\end{proof}

\begin{theorem}[KernelInductiveAcceptance obligation closure package]
\label{thm:kernel-inductive-acceptance-obligation-closure-package}
The carrier of \autoref{def:kernel-inductive-acceptance-carrier}, NameCert obligations of \autoref{thm:kernel-inductive-acceptance-namecert-obligations}, TasteGate nontrivial boundary of \autoref{thm:kernel-inductive-acceptance-nontrivial-boundary}, field-faithful boundary of \autoref{thm:kernel-inductive-acceptance-field-faithful-boundary}, constructor-ledger exhaustion of \autoref{thm:kernel-inductive-acceptance-constructor-ledger-exhaustion}, eliminator-ledger exhaustion of \autoref{thm:kernel-inductive-acceptance-eliminator-ledger-exhaustion}, and positivity-refusal nonescape of \autoref{thm:kernel-inductive-acceptance-refusal-nonescape} assemble the full paper-side obligation surface for $\KernelInductiveAcceptanceUp$. The assembled package exports no theorem-proof row, generic theorem trace, generator-closure theorem, type-membership theorem, subject-reduction target, or hidden kernel-correctness theorem.
\end{theorem}
\begin{proof}
Project the carrier tuple $I=(D,S,E,P,R,H,C,Q,N)$. These coordinates display the declaration candidate, constructor-signature ledger, emitted-eliminator ledger, positivity refusal, recursion discipline, transport, continuation, provenance, and local naming rows.

Apply \autoref{thm:kernel-inductive-acceptance-namecert-obligations}. Its five rows give carrier admission, constructor-signature exactness, emitted-eliminator exactness, positivity-refusal stability, and environment-entry non-escape for the same finite packet.

Use \autoref{thm:kernel-inductive-acceptance-nontrivial-boundary} and \autoref{thm:kernel-inductive-acceptance-field-faithful-boundary}. They enforce the AI TasteGate discipline by keeping the accepted and refused sides distinct and by making faithful reads preserve the paired constructor-signature plus emitted-eliminator fields.

The accepted read surface is exhausted by \autoref{thm:kernel-inductive-acceptance-constructor-ledger-exhaustion} and \autoref{thm:kernel-inductive-acceptance-eliminator-ledger-exhaustion}. Constructor consumers return to $S$, and recursor, induction-principle, and computation consumers pass through $S$ before $E$.

Finally apply \autoref{thm:kernel-inductive-acceptance-refusal-nonescape}. Failed positivity routes stay in $P$, so the package cannot leak into accepted declaration, trace, generator, type-classifier, or subject-reduction consumers through a hidden kernel-correctness theorem.
\end{proof}

\begin{theorem}[KernelInductiveAcceptance accepted declaration route]
\label{thm:kernel-inductive-acceptance-accepted-declaration-route}
Every accepted-declaration consumer of a $\KernelInductiveAcceptanceUp$ packet must pass through the displayed rows $D,S,E,R,H,C,Q,N$ of the carrier. The consumer may read the declaration candidate, constructor-signature ledger, emitted-eliminator ledger, recursion discipline, componentwise $\hsame$ transports, displayed $\Cont$ routes, $\Pkg$ provenance, and local $\NameCert$ rows, but it receives no hidden theorem-proof row and no accepted-declaration certificate outside those displayed rows.
\end{theorem}
\begin{proof}
Project the carrier of \autoref{def:kernel-inductive-acceptance-carrier}. Its accepted side is the finite route $D,S,E,R,H,C,Q,N$, with the refused positivity candidates separated into $P$.

Apply the NameCert obligations of \autoref{thm:kernel-inductive-acceptance-namecert-obligations}. Carrier admission, constructor-signature exactness, emitted-eliminator exactness, and environment-entry non-escape force an accepted-declaration read to expose $D$, then $S$, then $E$, with $R,H,C,Q,N$ carrying the route structure.

The field-faithful boundary of \autoref{thm:kernel-inductive-acceptance-field-faithful-boundary} prevents a consumer from dropping the signature or eliminator ledgers. The obligation package of \autoref{thm:kernel-inductive-acceptance-obligation-closure-package}, together with the existing TasteGate anchor recorded by this chapter, keeps the AI TasteGate read anchored in the displayed carrier rather than in a hidden theorem-proof row.
\end{proof}

\begin{theorem}[KernelInductiveAcceptance generator closure route]
\label{thm:kernel-inductive-acceptance-generator-closure-route}
Every generator-closure consumer of a $\KernelInductiveAcceptanceUp$ packet may read the constructor-signature ledger $S$ and emitted-eliminator ledger $E$ only through the displayed route rows $D,S,E,R,H,C,Q,N$ and the local $\NameCert_{\KernelInductiveAcceptanceUp}$ row. It cannot replace those rows by a generator-closure theorem, a theorem-proof trace, or an off-carrier closure certificate.
\end{theorem}
\begin{proof}
Begin with constructor-ledger and eliminator-ledger exhaustion, \autoref{thm:kernel-inductive-acceptance-constructor-ledger-exhaustion} and \autoref{thm:kernel-inductive-acceptance-eliminator-ledger-exhaustion}. They identify $S$ and $E$ as the only visible signature and generated-eliminator ledgers available to downstream generator consumers.

Use the NameCert surface to attach those ledgers to $D,R,H,C,Q,N$. The displayed $\Cont$ row records the route from declaration candidate to constructor signatures and from signatures to emitted eliminators, while $Q$ and $N$ supply provenance and local naming.

Finally apply the obligation closure package and the field-faithful boundary. A generator-closure consumer preserving the TasteGate surface must preserve the paired $S,E$ fields, so it cannot substitute a hidden theorem-proof trace or an off-carrier closure theorem for the displayed route.
\end{proof}

\begin{theorem}[KernelInductiveAcceptance subject reduction route]
\label{thm:kernel-inductive-acceptance-subject-reduction-route}
Every subject-reduction consumer of a $\KernelInductiveAcceptanceUp$ packet receives emitted eliminator and computation rows only after the constructor signatures in $S$ are visible. Failed positivity candidates remain in the refusal row $P$, and no subject-reduction consumer can transport such a candidate into the accepted route $D,S,E,R,H,C,Q,N$.
\end{theorem}
\begin{proof}
Start with the emitted-eliminator exactness clause in \autoref{thm:kernel-inductive-acceptance-namecert-obligations}. It requires every recursor, induction-principle, and computation read to pass through $E$ after the constructor-signature ledger $S$ is visible.

Apply \autoref{thm:kernel-inductive-acceptance-eliminator-ledger-exhaustion}. This pins the subject-reduction-facing rows to the displayed route $D,S,E,R,H,C,Q,N$ and rules out a bypass from declaration candidate directly to computation rows.

It remains to separate failed positivity candidates. By \autoref{thm:kernel-inductive-acceptance-refusal-nonescape}, such candidates land in $P$ and remain outside accepted declaration, trace, generator, type-classifier, and subject-reduction paths. The obligation closure package and the existing TasteGate anchor recorded by this chapter therefore leave subject reduction with only the displayed accepted route and no hidden theorem-proof row.
\end{proof}

\falsifiablePrediction{If the carrier is atomic, every accepted row exposes a constructor signature and emitted eliminator ledger; a packet lacking either row is rejected by the local classifier.}

\independenceWitness{kernelacceptancewitness, kernelacceptancetrace, generatorclosure, typeclassifiermembership, subjectreduction: none is bijective because those siblings lack the paired constructor-signature plus emitted-eliminator ledger row, while this carrier lacks their theorem-proof, generic trace, generator-closure, or subject-reduction target rows.}

\closureat{KernelInductiveAcceptanceUp}{obligationStr}

\begin{closurestatus}{\KernelInductiveAcceptanceUp}
  \constructivestory{A $\KernelInductiveAcceptanceUp$ packet is generated from an inductive declaration candidate, constructor-signature ledger, emitted-eliminator ledger, positivity-refusal row, recursion-discipline row, componentwise $\hsame$ transport, displayed $\Cont$ routes, $\Pkg$ provenance, and local naming.}
  \theoryclosure{\obligationClosure}
  \scopeclosed{\autoref{def:kernel-inductive-acceptance-carrier}, \autoref{thm:kernel-inductive-acceptance-namecert-obligations}, \autoref{thm:kernel-inductive-acceptance-nontrivial-boundary}, \autoref{thm:kernel-inductive-acceptance-field-faithful-boundary}, \autoref{thm:kernel-inductive-acceptance-constructor-ledger-exhaustion}, \autoref{thm:kernel-inductive-acceptance-eliminator-ledger-exhaustion}, \autoref{thm:kernel-inductive-acceptance-refusal-nonescape}, and \autoref{thm:kernel-inductive-acceptance-obligation-closure-package} bind the carrier, local NameCert obligation, nontrivial boundary, field-faithful boundary, constructor and eliminator ledger exhaustion, positivity refusal, environment nonescape, and obligation package rows.}
  \formalstatus{\auditCleanV}
  \leanchecked{BEDC.Derived.KernelInductiveAcceptanceUp.taste\_gate}
  \leantarget{BEDC.Derived.KernelInductiveAcceptanceUp.KernelInductiveAcceptanceTasteGate\_single\_carrier\_alignment}
  \bridgestatus{none}
  \origin{ai}
  \notclaimed{This chapter does not prove full kernel correctness, subject reduction, theorem-proof soundness, generic trace completeness, generic theorem proof rows, or any bridge theorem.}
  \upgradepath{Scoped closure requires binding the displayed kernel primitives and sibling consumer routes to one carrier-local readback surface.}
\end{closurestatus}
